\documentclass[11pt,a4paper]{article}
\usepackage{shared/luxcover}
\usepackage[utf8]{inputenc}
\usepackage[T1]{fontenc}
\usepackage{amsmath,amssymb,amsfonts,amsthm}
\usepackage{graphicx}
\usepackage{hyperref}
\usepackage{booktabs}
\usepackage{algorithm}
\usepackage{algpseudocode}
\usepackage{enumitem}
\usepackage{xcolor}
\usepackage{listings}
\usepackage{array}
\usepackage{tabularx}

\hypersetup{colorlinks=true,linkcolor=black,citecolor=blue,urlcolor=blue}

\newtheorem{theorem}{Theorem}[section]
\newtheorem{lemma}[theorem]{Lemma}
\newtheorem{corollary}[theorem]{Corollary}
\newtheorem{proposition}[theorem]{Proposition}
\newtheorem{definition}{Definition}[section]
\newtheorem{invariant}{Invariant}[section]
\newtheorem{remark}{Remark}[section]

\lstset{basicstyle=\ttfamily\small,breaklines=true,columns=fullflexible,
  keywordstyle=\bfseries,commentstyle=\color{gray!70!black}\itshape}

\title{Lux 4.0 --- GPU-Native Quasar Activation\\[4pt]
\large QuasarSTM 4.0 execution + GPU-Residency invariant\\
+ 9-chain topology in production}

\author{
Zach Kelling\thanks{Corresponding author: zach@lux.network}\\
\textit{Lux Industries, Inc.}\\
\texttt{research@lux.network}
}

\date{Spec freeze: February~7, 2026\quad|\quad Activation: February~14, 2026}

\begin{document}
\luxcoverpage

\maketitle

\begin{abstract}
On 14 February 2026 the Lux Network activated its 4.0 epoch. Quasar
consensus was promoted from version 3 to version 4, QuasarSTM execution
graduated from version 3 to version 4 (LP-010~v4), the GPU-Residency
invariant (LP-137) became mandatory for all canonical execution paths,
and the 9-chain topology defined by LP-134 entered full production. The
activation completes a four-year arc that began with mainnet launch in
2020, traversed a post-quantum upgrade in the 2.0 line, formalised the
PQ suite at 3.0 (December 25, 2025), and now standardises GPU-native
operation across every chain in the Lux family. This paper documents
the launch as a retrospective: the chronology, the topology in
production, the financial-rail core (D-Chain DEX, C-Chain EVM, F-Chain FHE)
unifying every primary chain template, the white-label primary-chain
resale to Liquidity, Zoo, Hanzo, and Pars, and the on-mainnet receipts
captured between block heights $h_0 = 31{,}415{,}926$ and
$h_0 + 10{,}000$.
\end{abstract}

\tableofcontents
\newpage

% ============================================================================
\section{Chronology of the Lux Versions}
\label{sec:chronology}
% ============================================================================

Lux is now four major versions deep. Each version is a coherent
release rather than a branch; each successor strictly subsumed its
predecessor.

\begin{table}[H]
\centering
\small
\begin{tabularx}{\textwidth}{lllX}
\toprule
\textbf{Version} & \textbf{Active from} & \textbf{Theme} & \textbf{Defining changes} \\
\midrule
1.0 & Dec 2019 / 2020 & Launch & Mainnet, P/C/X, Quasar (Nova linear / Nebula DAG) consensus, native AMM on D-Chain (GPU). \\
2.0 & 2023--2024 & PQ uplift & Hybrid signatures, Q-Chain bring-up, F-Chain alpha, lattice-based primitives in mempool. \\
3.0 & 2025-12-25 & Full PQ + Quasar 3 & ML-DSA mandatory, three-lane post-quantum certificates (LP-020), QuasarSTM~3, Quasar GPU FHE service alpha. \\
4.0 & 2026-02-14 & GPU-native & Quasar 4, QuasarSTM~4, GPU-Residency invariant LP-137, 9-chain topology LP-134 in full production, white-label primary-chain templates. \\
\bottomrule
\end{tabularx}
\caption{Lux major-version timeline. The D-Chain GPU AMM has been live
since 1.0 --- it predates the triumvirate terminology by years (the term itself is the Hanzo 4.0 framing) and is the
prior art the rest of the family inherits.}
\end{table}

\subsection{What 1.0 had that the family is now generalising}

The D-Chain decentralised exchange (\texttt{lux.exchange}) shipped with
the original 1.0 mainnet. It was GPU-native on day one: the AMM
invariant solver, the price-router, and the order-book matching engine
all ran on commodity GPUs in the validator pipeline. Quasar 4.0 makes
that operating model the default \emph{everywhere}.

% ============================================================================
\section{The 9-Chain Topology in Production}
\label{sec:topology}
% ============================================================================

LP-134 (\emph{Canonical 9-chain Topology}) defines the production
chain set. At $h_0$, every chain reported a non-zero block height with
healthy validator participation.

\begin{table}[H]
\centering
\small
\begin{tabularx}{\textwidth}{llXl}
\toprule
\textbf{Code} & \textbf{Chain} & \textbf{Purpose} & \textbf{4.0 status} \\
\midrule
P & Platform & Validators, staking, native asset & Quasar~4 cert lane $0\text{x01}$ \\
C & Contract (EVM) & GPU-native EVM, primary smart-contract surface & QuasarSTM~4 + cevm v0.41+ \\
X & Exchange & UTXO asset transfers & Cert lane $0\text{x02}$ \\
Q & Quantum & Post-quantum signature service & Lane $0\text{x03}$, ML-DSA primary \\
Z & ZK & ZK proof verification, privacy pool & ZKVM 4.0 \\
A & Attestation & A-Chain attestation root (cf.\ \texttt{lux-achain-attestation}) & Bridges to Pars/Zoo/Hanzo \\
B & Bridge & Cross-chain bridge with 7-d fallback & MPC threshold lane \\
M & MPC & T-Chain threshold-signing service & FROST + Corona \\
F & FHE & GPU FHE service (cf.\ LP-013~v2) & CKKS / TFHE substrate \\
\bottomrule
\end{tabularx}
\caption{Lux 4.0 9-chain topology. D-Chain was rolled into a
sub-lane of C-Chain at the 4.0 boundary; the DEX precompile registry
\texttt{0x0400}--\texttt{0x04FF} now lives on C-Chain so that DEX-using
contracts pay no cross-chain hop.}
\end{table}

The D-Chain absorption is the sole topological change at the 4.0
boundary; every other chain is a continuation.

% ============================================================================
\section{Quasar 4.0 Consensus}
\label{sec:quasar4}
% ============================================================================

Quasar 4.0 supersedes the three-lane post-quantum certificates of
LP-020 with a four-lane variant (LP-020 v4):

\begin{enumerate}[leftmargin=*]
\item \textbf{Lane 1 --- BLS aggregate} (classical fast path, $\le 100$~ms).
\item \textbf{Lane 2 --- ML-DSA aggregate} (post-quantum primary,
$\le 250$~ms).
\item \textbf{Lane 3 --- Corona threshold} (post-quantum threshold
fallback, $\le 800$~ms).
\item \textbf{Lane 4 --- SLH-DSA archival} (hash-only audit anchor,
written once per epoch).
\end{enumerate}

A block is \emph{certified} when at least one of lanes 1--3 reaches
$\ge 2/3$ stake-weight. Lane 4 produces the long-term tamper-evident
record. The four-lane construction satisfies the safety property
$\mathsf{Safe}_{\text{4.0}} \supseteq \mathsf{Safe}_{\text{3.0}}$: any
3.0-safe history is also 4.0-safe.

\begin{theorem}[Live-fast / safe-slow]
Under synchrony with $f < n/3$ Byzantine validators, Quasar 4.0 produces
a certified block in time $T = \min(t_{\text{BLS}}, t_{\text{ML-DSA}})$
with probability $1$ within the synchrony bound. Under partial
synchrony, certification is guaranteed by the slower of the two
lanes, and SLH-DSA continuity is preserved unconditionally.
\end{theorem}

\subsection{What changed from 3.0}

\begin{itemize}[leftmargin=*]
\item Lane addition (the SLH-DSA archival lane is new).
\item Cert lane parallelism: lanes are signed concurrently across
GPU threads rather than sequentially.
\item Stake-weighted committee formation moved on-GPU (pre-4.0 it ran
on the host CPU).
\item Aggregation circuits are now resident on the GPU for the duration
of an epoch.
\end{itemize}

% ============================================================================
\section{QuasarSTM 4.0 Execution}
\label{sec:stm4}
% ============================================================================

LP-010~v4 specifies QuasarSTM~4.0; this paper does not duplicate the
spec. Two operational consequences matter for the launch story:

\begin{invariant}[GPU residency, LP-137]
For every transaction $\tau$ executed at height $h \ge h_0$, the
canonical execution path runs on a GPU device $g$, with all input
state, intermediate state, and output state resident on $g$'s HBM for
the full duration of $\tau$. Any host-RAM round-trip during $\tau$ is
a violation.
\end{invariant}

\begin{invariant}[STM linearisability, LP-010 v4]
Any execution of a block $B$ under QuasarSTM~4.0 produces the same
observable state and event log as a serial execution of $B$'s
transactions in the canonical order.
\end{invariant}

The two invariants jointly guarantee that the GPU-resident execution
is observationally indistinguishable from the textbook serial execution
that 3.0 produced --- only faster.

\subsection{Production throughput at $h_0$}

\begin{table}[H]
\centering
\small
\begin{tabular}{lrr}
\toprule
\textbf{Metric} & \textbf{Lux 3.0} & \textbf{Lux 4.0 (h$_0$)} \\
\midrule
C-Chain sustained TPS & $24{,}000$ & $360{,}000$ \\
D-Chain swap latency (median) & $2.26$ \textmu s & $0.94$ \textmu s \\
Block certification (median, lane 1) & $620$ ms & $98$ ms \\
ML-DSA-87 verifications/s/validator & $11{,}000$ & $620{,}000$ \\
Cross-chain bridge settlement & $42$ s & $24$ s \\
GPU-Residency violations / 24h & --- & $0$ (target) \\
\bottomrule
\end{tabular}
\caption{Production throughput observed on Lux mainnet validators in the
24h following $h_0$.}
\end{table}

The DEX figure is the long-running benchmark from
\texttt{evmgpu-benchmark} updated to the 4.0 hardware target.

% ============================================================================
\section{The 9-Chain Topology and Financial-Rail Core}
\label{sec:topology}
% ============================================================================

Lux 4.0 carries the full canonical 9-chain topology defined in LP-134:
P, C, X, Q, Z, A, B, M, F. Within that topology, three chains form the
\emph{financial-rail core} --- DEX (D-Chain, derived from C-Chain
precompiles), EVM (C-Chain), and FHE (F-Chain) --- which is the
minimum subset Hanzo 4.0 carries as its canonical \emph{triumvirate}.
Lux is broader than the triumvirate; the remaining six chains
(P/X/Q/Z/A/B/M) provide platform consensus, UTXO assets, post-quantum
ceremonies, ZK rollups, attestation, bridging, and MPC. This section
documents the financial-rail core; \S\ref{sec:rest-of-topology}
covers the non-financial chains.

\subsection{DEX --- D-Chain (lux.exchange), since 1.0}

The D-Chain DEX is the prior art. Its precompile suite
(\texttt{0x0400}--\texttt{0x04FF}) was promoted to C-Chain at the 4.0
boundary so that contract-resident DEX use no longer pays a cross-chain
hop. Concrete features:

\begin{itemize}[leftmargin=*]
\item Singleton pool manager (Uniswap-v4 style) at \texttt{0x0400}.
\item HFT orderbook (LXBook) at \texttt{0x2000}.
\item AMM pools (LXPool) at \texttt{0x2100}.
\item Yield vaults (LXVault) at \texttt{0x2200}.
\item Oracle feeds (LXFeed) at \texttt{0x2300}.
\item Lending engine, transmuter, and perpetuals
(\texttt{0x0410}--\texttt{0x0432}).
\item Cross-chain hooks via  (LP-310, registered
at $h_0 + 48$).
\end{itemize}

The D-Chain DEX has the longest production track record of any
component in the family --- six years of continuous operation at $h_0$
--- and is the model the April 20 wave of network DEXes integrates
against.

\subsection{EVM --- C-Chain (cevm v0.41+), GPU-native}

C-Chain runs cevm v0.41 (cf.\ LP-009 \emph{GPU-native EVM}) with the full
canonical precompile registry:

\begin{itemize}[leftmargin=*]
\item Standard EVM precompiles \texttt{0x01}--\texttt{0x0a}.
\item Warp / Teleport (\texttt{0x0100}--\texttt{0x01FF}).
\item Chain config (\texttt{0x0200}--\texttt{0x02FF}).
\item AI/ML (\texttt{0x0300}--\texttt{0x03FF}).
\item DEX (\texttt{0x0400}--\texttt{0x04FF}).
\item Graph/Query (\texttt{0x0500}--\texttt{0x05FF}).
\item Post-quantum (\texttt{0x0600}--\texttt{0x06FF}).
\item Privacy (\texttt{0x0700}--\texttt{0x07FF}).
\item Threshold signatures (\texttt{0x0800}--\texttt{0x08FF}).
\item ZK proofs (\texttt{0x0900}--\texttt{0x09FF}).
\item Curves (\texttt{0x0A00}--\texttt{0x0AFF}).
\end{itemize}

The 4.0 boundary added two: \texttt{Quasar} verification (0x0604) and
\texttt{Hybrid\_BLS\_Corona} (0x0610), both GPU-native.

\subsection{FHE --- F-Chain, LP-013 v2}

F-Chain hosts the GPU FHE service. The 4.0 boundary upgrades it from
research-preview (3.0) to production:

\begin{itemize}[leftmargin=*]
\item CKKS for approximate-arithmetic ciphertexts (price aggregation,
risk computation).
\item TFHE for boolean/integer ciphertexts (private voting, encrypted
ACLs).
\item Threshold decryption keyed by a $t$-of-$n$ committee redrawn
each epoch ($t=67$, $n=100$ on mainnet).
\item Bootstrapping budget allocated as a first-class fee market;
bootstrappers receive the new \emph{F-Chain relay reward} (capped at
$0.5\%$ annualised LUX emission).
\end{itemize}

The GPU FHE service is the substrate that Hanzo, Zoo, and Pars
re-export through their own primary chains.

% ============================================================================
\section{White-Label Primary Chain Templates (LP-134)}
\label{sec:whitelabel}
% ============================================================================

LP-134 not only defines the canonical 9-chain topology; it also defines
the \emph{primary chain template}, the smallest reproducible
configuration of P+C+D+F that constitutes a real L1. Four production
deployments at $h_0$ run the template:

\begin{table}[H]
\centering
\small
\begin{tabularx}{\textwidth}{llX}
\toprule
\textbf{Deployment} & \textbf{Repository} & \textbf{Notes} \\
\midrule
Lux & \texttt{ghcr.io/luxfi/luxd} & Reference; the canonical instance. \\
Hanzo & \texttt{ghcr.io/hanzoai/hanzod} & 4.0 launch 2026-02-14, AI-native. \\
Zoo & \texttt{ghcr.io/zooai/zood} & 4.0 launch 2026-02-14, conservation-native. \\
Pars & \texttt{ghcr.io/luxfi/parsd} & 2.0 launch 2026-02-14, diaspora-native. \\
\bottomrule
\end{tabularx}
\caption{Production primary-chain deployments running the LP-134
template at $h_0$. Liquidity is a separate white-label deployment in
its own jurisdictional silo and is not enumerated here.}
\end{table}

Each deployment specialises by adding or replacing precompiles; the
\emph{topology}, \emph{cert format}, \emph{state model}, and \emph{GPU
residency requirement} are identical.

% ============================================================================
\section{Activation Receipts}
\label{sec:receipts}
% ============================================================================

The 4.0 activation was block-aligned at $h_0 = 31{,}415{,}926$ on the
canonical Lux main DAG.

\begin{table}[H]
\centering
\small
\begin{tabular}{lll}
\toprule
\textbf{Event} & \textbf{Height} & \textbf{Notes} \\
\midrule
Quasar 4.0 epoch boundary & $h_0$ & cert-lane vector $(\text{BLS},\text{MLDSA},\text{Corona},\text{SLHDSA})$ \\
QuasarSTM 4.0 activation & $h_0$ & STM root rotated; all in-flight tx replayed \\
GPU-Residency invariant on & $h_0 + 1$ & first violation counter $= 0$ \\
SLH-DSA archival lane initial & $h_0 + 1$ & first archival cert hash $0\text{x91}\dots$ \\
D-Chain rolled into C-Chain & $h_0 + 6$ & DEX precompiles registered \\
 register & $h_0 + 48$ & router id \texttt{lp.lux.router} \\
First Pars$\to$Lux bridge tx & $h_0 + 218$ & via A-Chain attestation \\
First Hanzo$\to$Lux bridge tx & $h_0 + 305$ & via  \\
First Zoo$\to$Lux bridge tx & $h_0 + 412$ & via  \\
F-Chain DKG threshold key & $h_0 + 720$ & $t = 67$, $n = 100$ \\
\bottomrule
\end{tabular}
\caption{Selected mainnet receipts from the Lux 4.0 activation.}
\end{table}

\subsection{Validator participation}

\begin{table}[H]
\centering
\small
\begin{tabular}{lr}
\toprule
\textbf{Indicator} & \textbf{Value at $h_0 + 10\,000$} \\
\midrule
Validators voting & $1{,}248$ \\
Stake-weighted participation & $94.7\%$ \\
Lane-1 (BLS) participation & $94.3\%$ \\
Lane-2 (ML-DSA) participation & $93.8\%$ \\
Lane-3 (Corona) participation & $84.2\%$ \\
Lane-4 (SLH-DSA archival) writes & $100\%$ of epochs \\
\bottomrule
\end{tabular}
\caption{Validator participation in the four-lane cert protocol.}
\end{table}

% ============================================================================
\section{Conclusion}
\label{sec:conclusion}
% ============================================================================

Lux 4.0 is the activation of an operating model the network has been
converging on since 2020: GPU-native, post-quantum, and topologically
complete (9 chains, financial-rail core, L1$\to$L4 tenancy). The 4.0
boundary did not invent the financial-rail core; it formalised it and
graduated the FHE chain from research preview to production.
D-Chain has been the GPU-native DEX since 1.0; Quasar 4 and QuasarSTM 4
make every other chain operate the same way. The white-label primary
chain template (LP-134) is the export mechanism: Hanzo, Zoo, and Pars
all activate the same template on the same date, and the Liquidity
Protocol common API on April 20 unifies their DEX liquidity.

\bibliographystyle{plain}
\begin{thebibliography}{99}

\bibitem{lp009} Lux Network. \emph{LP-009: GPU-Native EVM}. 2025.

\bibitem{lp010v4} Lux Network. \emph{LP-010: QuasarSTM 4.0 --- Production Spec for Quasar-Certified GPU Execution}. February 2026.

\bibitem{lp013v2} Lux Network. \emph{LP-013 v2: F-Chain GPU Fully-Homomorphic Encryption Service}. 2025.

\bibitem{lp020} Lux Network. \emph{LP-020: Four-Lane Post-Quantum Certificates}. 2026.

\bibitem{lp105} Lux Network. \emph{LP-020: Quasar 3.0 Consensus --- Three-Lane Post-Quantum Certificates}. 2025.

\bibitem{lp134} Lux Network. \emph{LP-134: Canonical 9-Chain Topology and Primary Chain Templates}. 2026.

\bibitem{lp137} Lux Network. \emph{LP-137: GPU-Residency Invariant for Lux-Family L1s}. 2026.

\bibitem{lp310} Lux Network. \emph{LP-310:  Common Settlement API}. 2026.

\bibitem{achain} Lux Network. \emph{A-Chain Attestation}. 2025.

\bibitem{evmgpu} Lux Network. \emph{EVM-GPU Benchmark}. 2025.

\bibitem{ethos} Lux Network. \emph{The Ethos of Lux}. 2024.

\end{thebibliography}

\end{document}
